\subsection{Additive Functions of Modules}

Let A be a ring, and let C be a set of classes of A-modules (VIII, p.~51); we say that an A-module is of type C if its class belongs to C.

DEFINITION 1. --- We say that the set C of classes of A-modules is additive if every zero module is of type C and the direct sum of two modules of type C is of type C. We say that C is hereditary if it is additive and the submodules and quotient modules of a module of type C are of type C.

Examples. --- 1) The set of classes of A-modules of finite length is hereditary (II, §1, No.~10, p.~212, Proposition 16).

\begin{enumerate}
\def\labelenumi{\arabic{enumi})}
\setcounter{enumi}{1}
\item
  The set of classes of finitely generated A-modules is additive. If the ring A is Noetherian, then this set is hereditary (VIII, p.~3, Proposition 3 and VIII, p.~7, Proposition 4).
\item
  The set of classes of finitely generated projective A-modules is additive but generally not hereditary.
\end{enumerate}

DEFINITION 2. --- Let \(\varphi\) be a mapping from \(\mathcal{C}\) to an abelian group G (written additively); set \(\varphi(\mathrm{E}) = \varphi(\mathrm{cl}(\mathrm{E}))\) for every A-module E of type \(\mathcal{C}\) . We say that \(\varphi\) is an additive function of modules (resp. a weakly additive function of modules) if we have \(\varphi(\mathrm{E}) = \varphi(\mathrm{E}') + \varphi(\mathrm{E}'')\) for every exact sequence (resp. for every split exact sequence)

\[
0 \longrightarrow \mathrm{E} ^ {\prime} \longrightarrow \mathrm{E} \longrightarrow \mathrm{E} ^ {\prime \prime} \longrightarrow 0
\]

of modules of type C.

Example 4. --- Let C be the set of classes of A-modules of finite length. The mapping \(\operatorname{long}_{\mathbf{A}} : \mathcal{C} \to \mathbf{Z}\) that sends a class of A-modules of finite length to its length is an additive function of modules (II, §1, No.~10, p.~213, Corollary 3). The results of this subsection are a generalization of the results on modules of finite length established in II, §1, No.~10, p.~212--214.

In the remainder of this subsection, we consider an additive set C of A-modules and an additive mapping \(\varphi\) from C to an abelian group G.

Let \(E\) and \(E'\) be modules of type \(\mathcal{C}\) ; then \(E \oplus E'\) is of type \(\mathcal{C}\) , and there exists a split exact sequence (II, §1, No.~9, p.~210)

\[
0 \longrightarrow \mathrm{E} \longrightarrow \mathrm{E} \oplus \mathrm{E} ^ {\prime} \longrightarrow \mathrm{E} ^ {\prime} \longrightarrow 0;
\]

from this, we deduce

\[
\varphi (\mathrm{E} \oplus \mathrm{E} ^ {\prime}) = \varphi (\mathrm{E}) + \varphi (\mathrm{E} ^ {\prime}).\tag{1}
\]

In particular, we have \(\varphi(0) = 0\) .

PROPOSITION 1. --- Suppose that \(\mathcal{C}\) is hereditary. Let E and F be A-modules and \(u: \mathrm{E} \to \mathrm{F}\) a linear mapping.

\begin{enumerate}
\def\labelenumi{\alph{enumi})}
\item
  If E or F is of type C, then so is the image of u.
\item
  If E is of type C, then so is the kernel of u, and we have
\end{enumerate}

\[
\varphi (\mathrm{E}) = \varphi (\mathrm{Ker} u) + \varphi (\mathrm{Im} u).\tag{2}
\]

\begin{enumerate}
\def\labelenumi{\alph{enumi})}
\setcounter{enumi}{2}
\tightlist
\item
  If \(\mathrm{F}\) is of type \(\mathcal{C}\) , then so is the cokernel of \(u\) , and we have
\end{enumerate}

\[
\varphi (\mathrm{F}) = \varphi (\operatorname{Im} u) + \varphi (\operatorname{Coker} u).\tag{3}
\]

The proposition follows from the existence of exact sequences

\[
0 \longrightarrow \operatorname{Ker} u \longrightarrow \mathrm{E} \longrightarrow \operatorname{Im} u \longrightarrow 0,
\]

\[
0 \longrightarrow \operatorname{Im} u \longrightarrow \mathrm{F} \longrightarrow \operatorname{Coker} u \longrightarrow 0.
\]

COROLLARY. --- Let \((\mathrm{E}_{i})_{0\leqslant i\leqslant n}\) be a finite sequence of modules of type C. If there exists an exact sequence

\[
0 \longrightarrow \mathrm{E} _ {0} \stackrel {u _ {0}} {\longrightarrow} \mathrm{E} _ {1} \stackrel {u _ {1}} {\longrightarrow} \dots \longrightarrow \mathrm{E} _ {n - 1} \stackrel {u _ {n - 1}} {\longrightarrow} \mathrm{E} _ {n} \longrightarrow 0,
\]

then we have

\[
\sum_ {i = 0} ^ {n} (- 1) ^ {i} \varphi (\mathrm{E} _ {i}) = 0.\tag{4}
\]

Let us prove the corollary by induction on \(n\) ; the cases \(n = 0\) and \(n = 1\) are trivial. So let \(n \geqslant 1\) , and let

\[
0 \longrightarrow \mathrm{E} _ {0} \stackrel {u _ {0}} {\longrightarrow} \mathrm{E} _ {1} \stackrel {u _ {1}} {\longrightarrow} \dots \longrightarrow \mathrm{E} _ {n - 1} \stackrel {u _ {n - 1}} {\longrightarrow} \mathrm{E} _ {n} \stackrel {u _ {n}} {\longrightarrow} \mathrm{E} _ {n + 1} \longrightarrow 0
\]

be an exact sequence of modules of type \(\mathcal{C}\) . By Proposition 1, the kernel \(F\) of \(u_{n}\) is a module of type \(\mathcal{C}\) , and we have

\[
\varphi (\mathrm{F}) = \varphi (\mathrm{E} _ {n}) - \varphi (\mathrm{E} _ {n + 1}).\tag{5}
\]

Moreover, we have an exact sequence

\[
0 \longrightarrow \mathrm{E} _ {0} \longrightarrow \mathrm{E} _ {1} \longrightarrow \mathrm{E} _ {2} \longrightarrow \dots \longrightarrow \mathrm{E} _ {n - 1} \longrightarrow \mathrm{F} \longrightarrow 0,
\]

and the induction hypothesis gives the relation

\[
\sum_ {i = 0} ^ {n - 1} (- 1) ^ {i} \varphi (\mathrm{E} _ {i}) + (- 1) ^ {n} \varphi (\mathrm{F}) = 0.\tag{6}
\]

We then immediately deduce from (5) and (6) that \(\sum_{i=0}^{n+1}(-1)^{i}\varphi(\mathrm{E}_{i}) = 0\) , and the corollary follows.

PROPOSITION 2. --- Suppose that the set C is hereditary. Let E be an A-module, and let M and N be submodules of E.

\begin{enumerate}
\def\labelenumi{\alph{enumi})}
\tightlist
\item
  If the modules M and N are of type C, then so are the modules \(M \cap N\) and \(M + N\) , and we have
\end{enumerate}

\[
\varphi (\mathrm{M} + \mathrm{N}) + \varphi (\mathrm{M} \cap \mathrm{N}) = \varphi (\mathrm{M}) + \varphi (\mathrm{N}).
\]

\begin{enumerate}
\def\labelenumi{\alph{enumi})}
\setcounter{enumi}{1}
\tightlist
\item
  If the modules E/M and E/N are of type C, then so are the modules \(\mathrm{E}/(\mathrm{M}\cap\mathrm{N})\) and \(\mathrm{E}/(\mathrm{M}+\mathrm{N})\) , and we have
\end{enumerate}

\[
\varphi (\mathrm{E} / (\mathrm{M} + \mathrm{N})) + \varphi (\mathrm{E} / (\mathrm{M} \cap \mathrm{N})) = \varphi (\mathrm{E} / \mathrm{M}) + \varphi (\mathrm{E} / \mathrm{N}).
\]

Assertions a) and b) follow from the existence of the exact sequences

\[
0 \to \mathrm{M} \cap \mathrm{N} \to \mathrm{M} \oplus \mathrm{N} \to \mathrm{M} + \mathrm{N} \to 0
\]

and

\[
0 \to \mathrm{E} / (\mathrm{M} \cap \mathrm{N}) \to (\mathrm{E} / \mathrm{M}) \oplus (\mathrm{E} / \mathrm{N}) \to \mathrm{E} / (\mathrm{M} + \mathrm{N}) \to 0
\]

(II, §1, No.~7, p.~207, Proposition 10).

PROPOSITION 3. --- Suppose that \(\mathcal{C}\) is hereditary. Let E be a module of type \(\mathcal{C}\) and \((\mathrm{E}_i)_{0\leqslant i\leqslant n}\) a composition series of E (I, §4, No.~7, p.~41). For \(1\leqslant i\leqslant n\) , the module \(\mathrm{E}_{i - 1} / \mathrm{E}_i\) is of type \(\mathcal{C}\) , and we have

\[
\varphi (\mathrm{E}) = \sum_ {i = 1} ^ {n} \varphi (\mathrm{E} _ {i - 1} / \mathrm{E} _ {i}).
\]

Since C is hereditary, the modules \(E_{0}/E_{1}\) and \(E_{1}\) are of type C, and we have \(\varphi(\mathrm{E}) = \varphi(\mathrm{E}_{0}/\mathrm{E}_{1}) + \varphi(\mathrm{E}_{1})\) . Since the sequence \((\mathrm{E}_{i+1})_{0 \leqslant i \leqslant n-1}\) is a composition series of \(E_{1}\) , Proposition 3 follows by induction on n.

